\scanpage{61}
\mspage{319}

Let us now consider the homomorphism \unsure{(3)}
\[
\tag{18}
\left\{
\begin{aligned}
\widehat\varphi:\quad &\widehat\Pi_{0,3}\longrightarrow
  \widehat{\mathfrak S}_{0,3}^{+}\\
&\ell_0\longmapsto\ell'_0=\sigma_\infty\rho=\varepsilon_0\\
&\ell_1\longrightarrow\ell'_1=\sigma_\infty\\
&\ell_\infty\longmapsto\ell'_\infty=\rho^{-1}
\end{aligned}
\right.
\]
and let us write that it commutes, \unsure{as} an outer homomorphism, with $u$.
This therefore means that there exists a
\[
\tag{19}
\lambda\in\widehat{\mathfrak S}_{0,3}^{+}
\]
such that we have
\[
\tag{20}
\widehat\varphi(\unsure{ux})=\lambda\,u\widehat\varphi(x)\,\lambda^{-1}
\qquad \forall x\in\widehat\Pi_{0,3},
\]
and it suffices to express this condition for $x$ equal to one of the two
generators $\ell_0,\ell_1$ of $\widehat\Pi_{0,3}$

For $x=\ell_0$ relation (20) is written
\[
\tag{21}
u(\ell'_0)=\operatorname{int}(\lambda^{-1}\widehat\varphi(g_0))(\ell_0'{}^{p})
\]
which, compared with (17), gives
\[
\operatorname{int}(g_0^{-1}\lambda^{-1}\widehat\varphi(g_0))(\ell_0'{}^{p})
=\ell_0'{}^{p}\quad\text{hence}
\]
\[
\deleted{(22)}\qquad
g_0^{-1}\lambda^{-1}\widehat\varphi(g_0)=\ell_0'{}^{\beta}
\qquad(\beta\in\widehat{\mathbb Z})
\]
or again
\[
\tag{22}
\boxed{\lambda=\widehat\varphi(g_0)\ell_0'{}^{-\beta}g_0^{-1}}
\inserted{=\widehat\varphi(g_0h_0^{-1})=\widehat\varphi(\sigma_\infty(g_0))}
\]
\margin{$\searrow$ \unsure{cf.} below (24 bis)}

For $x=\ell_1$ relation (20) is written

\scanpage{62}

\[
\tag{23}
u(\ell'_1)=\operatorname{int}(\lambda^{-1}\widehat\varphi(g_1))(\ell_1'{}^{p})
\]
\margin{where $(g_1=\sigma_\infty(g_0)=\sigma_\infty g_0\sigma_\infty\unsure{)}$}
or again
\[
\sigma_\infty=\operatorname{int}(\lambda^{-1}\widehat\varphi(g_1))(\sigma_\infty)
\]
\margin{For $\sigma_\infty^{p}=\sigma_\infty$ since
\inserted{$\sigma_\infty^2=1$ and} \unsure{the image} of $p$ in
$(\mathbb Z/2\mathbb Z)^*$ is $1$!}
or again $\sigma_\infty$ commutes with
\[
\lambda^{-1}\widehat\varphi(g_1)
\overset{\inserted{\text{cf. (22)}}}{=}
g_0\ell_0'{}^{\beta}\widehat\varphi
 \bigl(\underbrace{g_0^{-1}\sigma_\infty(g_0)}_{h_0^{-1}}\bigr)
=g_0\ell_0'{}^{\beta}\widehat\varphi(h_0^{-1})
\]
or again with its inverse $\widehat\varphi(h_0)\ell_0'{}^{-\beta}g_0^{-1}$
which is also written
\[
\underbrace{\sigma_\infty
 \bigl(\widehat\varphi(h_0).\ell_0'{}^{-\beta}.g_0^{-1}\bigr)}
=\widehat\varphi(h_0).\ell_0'{}^{-\beta}.g_0^{-1}
\]
\[
\sigma_\infty(\widehat\varphi(h_0)).
\underbrace{\varepsilon_1^{-\beta}}_{(\rho\sigma_\infty)^{-\beta}}.
\sigma_\infty(g_0)^{-1}
=\widehat\varphi(h_0)\ell_0'{}^{-\beta}g_0^{-1}
\]
\[
\sigma_\infty(\widehat\varphi(h_0))
 (\underbrace{\rho\sigma_\infty}_{\varepsilon_1})^{-\beta}h_0
=\widehat\varphi(h_0).\varepsilon_0^{-\beta}
\]
\[
\tag{24}
\boxed{
\overbrace{\varepsilon_1^{-\beta}h_0\varepsilon_0^{\beta}}
 ^{\inserted{\sigma_\infty(g_0\varepsilon_0^{\beta})^{-1}(g_0\varepsilon_0^{\beta})}}
=\sigma_\infty(\widehat\varphi(h_0))^{-1}\widehat\varphi(h_0)
}
\]
\margin{\unsure{Correct?} $\searrow$}
\margin{\uline{N.B.} If $g_0$ is replaced by
\inserted{$g'_0=$} $g_0\ell_0^{\lambda}$ $(\lambda\in\widehat{\mathbb Z})$,
$h_0$ is replaced by $h'_0=\ell_1^{-\lambda}\unsure{h_0|g_0}\ell_0^{\lambda}$,
and we \unsure{find} that $\beta$ \unsure{is} \unsure{''} replaced by
$\beta'=\beta-3\lambda$.
This shows \inserted{?} that we can in a \uline{unique way}
\illegible{1} \uline{normalize} the choice of $g_0$, in such a way that we have
$\beta\in\{0,1,2\}$.
(It would be necessary to elucidate whether $\beta$ actually takes all
\unsure{three} values modulo $3$\ldots)}
\margin{\uuline{N.B.} This equation is equivalent to the
\uline{simpler} equation
\[
\text{(24 bis)}\qquad
\boxed{\widehat\varphi(h_0)=\sigma_\infty^{\unsure{\varepsilon}}g_0\varepsilon_0^{\beta}}
\]}
\[
\unsure{(25)}\qquad
\left\{
\begin{aligned}
&h_0\text{ given by (8) in terms of }g_0\\
&\varepsilon_0=\ell'_0=\sigma_\infty\rho
   &&(\varepsilon_0^2=\ell_0)\\
&\varepsilon_1=\rho\sigma_\infty=\sigma_0\rho
   &&(\varepsilon_1^2=\ell_1)
\end{aligned}
\right.
\]

Setting, for $\deleted{\unsure{g}}\ \mu\in
\widehat{\mathfrak S}_{0,3}^{+}$,
\[
\tag{26}
\widetilde\mu=\sigma_\infty(\mu)^{-1}\mu
=\sigma_\infty\mu^{-1}\sigma_\infty\mu=[\sigma_\infty,\mu]
\]
(so that $h_0=\widetilde g_0$) the formula \unsure{(24)|(25)} takes the form
\[
\tag{27}
\varepsilon_1^{-\beta}\unsure{\widetilde g_0}\varepsilon_0^\beta
=\bigl(\widehat\varphi(\widetilde g_0)\bigr)^{\sim}.
\]

\scanpage{63}
\mspage{320}

In summary:

\uline{Proposition} Let us consider an outer automorphism $u$ of the group with loops
$\widehat\Pi_{0,3}$, realized by $g_0\in\widehat\Pi_{0,3}$
\inserted{modulo $L_0$,}
and $p\in\widehat{\mathbb Z}^{*}$,
(and $\gamma\in\widehat{\mathbb Z}^{*}$ uniquely determined
\unsure{by terms|in terms} of $g_0$ and $p$)
by formulas (11), by means of an actual automorphism
\inserted{(still denoted $u$)} which commutes with $\sigma_\infty$ ---
the data $g_0,p$ (and $\gamma$) \unsure{being} \inserted{therefore}
required to satisfy conditions (9) and (10) above,
plus the generation condition (11 bis).
\inserted{(We extend $u$ \unsure{to} an automorphism of
$\widehat{\mathfrak S}_{0,3}^{+}$
by $\operatorname{int}(u)$, cf. (15))}
For the outer homomorphism
$\widehat\varphi:\widehat\Pi_{0,3}\longrightarrow
\widehat{\mathfrak S}_{0,3}^{+}$
defined by (18) to commute with the outer action of $u$,
it is necessary and sufficient that we can find a $\beta\in\widehat{\mathbb Z}$,
such that we have (24), or what amounts to the same, (27),
\margin{\unsure{better:} \uuline{24 bis}}

This condition \inserted{(for $\beta$ fixed)} is also equivalent to formula
(20) for $x=\ell_\infty$ ($\lambda$ \unsure{being} always given by (22)),
\deleted{\unsure{(2}} i.e. to the formula
\[
\tag{28}
u(\ell'_\infty)=\operatorname{int}(\lambda^{-1}\widehat\varphi(g_\infty))
 (\ell_\infty'{}^p),\qquad\text{where}
\]
\[
\tag{29}
g_\infty=g_0\ell_0^{-\gamma}\rho^{-1}(h_0)
\qquad(\text{and }\lambda\text{ given by (22)}).
\]

\scanpage{64}

formula which thus makes precise the \unsure{presumptive} formula (18). It is also written, more explicitly
\[
\rho^{-1}g^{-1}
=\operatorname{int}\bigl(g_0\varepsilon_0^{\beta}
\widehat\varphi(\ell_0^{\gamma}\rho^{-1}(h_0))\bigr)
(\rho^{-p})
\]
\inserted{\[
g_0\varepsilon_0^{\beta+\gamma}\widehat\varphi(\rho^{-1}(h_0))
\]}
i.e. (taking inverses, and making $g$ explicit)
\[
(30)\qquad
\boxed{
\deleted{\illegible{1}}\inserted{h_0^{-1}\ell_1^{\gamma}\rho}
=\operatorname{int}\bigl(
\deleted{\unsure{g_0}}\varepsilon_0^{\deleted{\unsure{\beta+\gamma}}\inserted{\alpha}}
\widehat\varphi(\rho^{-1}(h_0))\bigr)
(\rho^{p\deleted{\unsure{-p}}})
}
\]
\margin{\uline{N.B.} We put $\alpha=\beta+\gamma$. The existence \unsure{of} $\beta\in\widehat{\mathbb Z}$ satisfying (24), is equivalent to that \unsure{of} $\alpha\in\widehat{\mathbb Z}$ satisfying (30).}

The equivalence of (24) and (30) may seem a little strange, owing to the fact \inserted{for example} that in (24) the exponent $p$ does not enter explicitly, but only $h_0$, whereas formula (30) brings in the value of $p\in\widehat{\mathbb Z}^{*}$ modulo $3$ (which gives $\rho^p=\rho$ or $\rho^p=\rho^{-1}$). But in any event, whether it be in the form (24) \unsure{where|or} in the form (30), it does seem that the condition in question is highly non-tautological!

To do things properly, it would be necessary to \unsure{first} make explicit the congruential conditions of \S~39 in terms of $g_0,p$.

\scanpage{65}
\mspage{321}

I would like to \unsure{compare} the outer automorphisms of $\widehat\Pi_{0,3}$, and those of $\widehat{\mathfrak S}_{0,3}^{+}$, taking account of the exact sequence
\[
(31)\qquad
\diagram{\begin{tikzcd}[ampersand replacement=\&]
1\arrow[r,no head]\&\widehat\Pi_{0,3}\arrow[r,no head]\&
\widehat{\mathfrak S}_{0,3}^{+}\arrow[r]\arrow[d,draw=none,"{\inserted{\unsure{\wr}}}" description]\&
\mathfrak S_{3}\arrow[r,no head]\&1\\
\&\&\inserted{\widehat{\Pi'}_{0,3}}
\end{tikzcd}}
\]

\deleted{\uline{Proposition}} \inserted{Let us note} that every outer automorphism \inserted{$u$} of $\widehat{\mathfrak S}_{0,3}^{+}$ which \unsure{leaves invariant} the normal subgroup $\widehat\Pi_{0,3}$, and induces on the quotient $\mathfrak S_{3}$ a trivial outer automorphism, can be written as obtained from an automorphism \inserted{$u_0$} of $\widehat{\mathfrak S}_{0,3}^{+}$, \unsure{stabilizing} $\widehat\Pi_{0,3}$, and inducing the \uline{identity} on the quotient $\mathfrak S_{3}$ --- and $u_0$ is then determined modulo composition with an element $\operatorname{int}(g)$, where $g\in\widehat{\mathfrak S}_{0,3}^{+}$ \unsure{is} such that its image in $\mathfrak S_{3}$ is central, i.e. $=1$, i.e. $g\in\widehat\Pi_{0,3}$. Thus
\[
(32)\qquad
\begin{gathered}
\text{group of outer automorphisms}\\
\text{of $\widehat{\mathfrak S}_{0,3}^{+}$ which}\\
\text{induce on $\mathfrak S_{3}$ a}\\
\text{trivial outer automorphism (hence \unsure{leave invariant}}\\
\widehat\Pi_{0,3})
\end{gathered}
\simeq
\begin{gathered}
\text{group of}\\
\text{automorphisms of}\\
\widehat{\mathfrak S}_{0,3}^{+}\text{ which}\\
\text{induce}\\
\text{the identity on}\\
\text{its quotient $\mathfrak S_{3}$,}\\
\text{\uline{modulo} the sub-}\\
\text{group of}\\
\text{inner automorphisms}\\
\text{by $g\in\widehat\Pi_{0,3}$}
\end{gathered}
\]